% !TEX program = xelatex
%%% ---------------------------------------------------------------------------
%%% main.tex —— 试卷入口（默认是空骨架：只有抬头、评分表和各栏占位，不含任何题目）
%%%
%%% 出卷步骤：
%%%   1. 把下面所有「〈…〉」占位内容换成自己的信息；
%%%   2. 按注释里的用法在各栏里写题目；
%%%   3. 编译（XeLaTeX，连续两遍；Overleaf：Settings → Compiler 选 XeLaTeX）。
%%%
%%% 想看完整示例（虚构题目 + 插图 + 公式）？把下面的
%%%   \showexamplefalse   改成   \showexampletrue   再编译即可。
%%% ---------------------------------------------------------------------------
\documentclass[12pt,a4paper]{exampaper}
%% 字体集默认自动检测（Windows→windows，macOS→mac，Linux→fandol）。
%% 需要指定时写成 \documentclass[fontset=fandol]{exampaper}

%% 页脚左侧的试卷名（页脚显示：试卷名  共 X 页  第 Y 页）
\examname{〈试卷名〉}

\newif\ifshowexample
\showexamplefalse          % ← 改成 \showexampletrue 看完整示例

\begin{document}

\ifshowexample

  %% ===================== 完整示例（虚构内容）=====================
  %%% ---------------------------------------------------------------------------
%%% example-exam.tex —— 完整示例（虚构内容，仅演示排版与用法）
%%%
%%% 本文件是「内容文件」，不含 \documentclass，由 main.tex 在
%%% \showexampletrue 时 \input 进来。想看效果：把 main.tex 里的
%%% \showexamplefalse 改成 \showexampletrue 再编译。
%%%
%%% 演示内容：抬头与信息栏、评分表、填空题、选择题、计算题（答题留白）、
%%% 应用题（插图 + 公式），以及页脚与左侧密封线。
%%% ---------------------------------------------------------------------------

\examname{示例课程期末试卷(A)}          % 页脚左侧的试卷名

%% ======================= 抬头 =======================
\examtitle{示例大学试卷（A 卷）}

\begin{center}
  \examfield[3.4cm]{课程名称}{示例课程}\hspace{0.4em}%
  \examfield[2.2cm]{课程代码}{EXM1001}\hspace{0.4em}%
  \examfield[1.8cm]{考试学期}{25-26-1}\\[0.4em]
  \examfield[4.4cm]{适用专业}{示例专业}\hspace{0.4em}%
  \examfield[2.0cm]{考试形式}{闭\hspace{0.5em}卷}\hspace{0.4em}%
  \examfield[1.9cm]{考试时长}{120\hspace{0.4em}分钟}
\end{center}

\vspace{0.8em}

\examscoretable[4]                      % 4 个大题

\vspace{0.6em}

%% ======================= 一、填空题 =======================
\examsection{填空题（本题共 3 小题，每小题 4 分，满分 12 分）}
\begin{examquestions}
  \examquestion 若 $f(x)=x^{2}$，则 $f'(1)=$ \underline{\hspace{2.5cm}}。
  \examquestion 设 $D=[0,1]\times[0,1]$，则 $\iint_{D} x\,\mathrm{d}\sigma=$
    \underline{\hspace{2.5cm}}。
  \examquestion 设矩阵 $A=\begin{pmatrix}1&2\\3&4\end{pmatrix}$，
    则 $\det A=$ \underline{\hspace{2.5cm}}。
\end{examquestions}

%% ======================= 二、选择题 =======================
\examsection{单项选择题（本题共 2 小题，每小题 5 分，满分 10 分）}
\begin{examquestions}
  \examquestion 函数 $y=\sin x$ 的最小正周期是
    \quad A.\ $\pi$ \quad B.\ $2\pi$ \quad C.\ $4\pi$ \quad D.\ $\dfrac{\pi}{2}$
  \examquestion 下列级数中收敛的是
    \quad A.\ $\sum_{n=1}^{\infty}\dfrac{1}{n}$ \quad
    B.\ $\sum_{n=1}^{\infty}\dfrac{1}{n^{2}}$\quad
    C.\ $\sum_{n=1}^{\infty}(-1)^{n}$ \quad D.\ $\sum_{n=1}^{\infty}2^{n}$
\end{examquestions}

%% ======================= 三、计算题（小题之间留答题空间）=======================
\examsection{计算题（本题共 2 小题，每小题 8 分，满分 16 分）}
\begin{examquestions}[7cm]              % 小题之间留 7cm 答题空间
  \examquestion 计算 $\displaystyle\lim_{x\to 0}\frac{\sin x}{x}$。
  \examquestion 计算 $\displaystyle\int_{0}^{1} x^{2}\,\mathrm{d}x$。
\end{examquestions}

%% ======================= 四、应用题（含插图）=======================
\examsection{应用题（本题共 1 小题，满分 12 分）}
\begin{examproblem}
  \examquestion 设某物理量的变化曲线如图~\ref{fig:curve} 所示，
  该曲线在 $t=0$ 附近经过原点。试估计曲线在 $t=1$ 处的函数值，
  并说明你的估算依据。

  %% 插图：\examfigure[<宽度>]{<图题>}{<图片文件>}
  %% 图片放在项目根目录的 figures/ 下；png / jpg / pdf 均可
  \examfigure[0.62\textwidth]{示例曲线（演示插图用法）}{figures/sample-figure.png}
  \label{fig:curve}

  %% 需要更自由的图文排版时，直接用标准 figure 环境（graphicx 已加载）：
  %% \begin{figure}[htbp]
  %%   \centering
  %%   \includegraphics[width=0.5\textwidth]{figures/your-image.png}
  %%   \caption{图题}\label{fig:your}
  %% \end{figure}
\end{examproblem}


\else

  %% ===================== 空白骨架 =====================
  \examtitle{〈学校名〉试卷（〈A 卷〉）}

  %% 信息栏：\examfield[<下划线宽度>]{<栏目标题>}{<填写内容>}
  \begin{center}
    \examfield[3.4cm]{课程名称}{〈课程名称〉}\hspace{0.4em}%
    \examfield[2.2cm]{课程代码}{〈课程代码〉}\hspace{0.4em}%
    \examfield[1.8cm]{考试学期}{〈20xx-xx-x〉}\\[0.4em]
    \examfield[4.4cm]{适用专业}{〈适用专业〉}\hspace{0.4em}%
    \examfield[2.0cm]{考试形式}{闭\hspace{0.5em}卷}\hspace{0.4em}%
    \examfield[1.9cm]{考试时长}{〈150〉\hspace{0.4em}分钟}
  \end{center}

  \vspace{0.8em}

  %% 评分表：\examscoretable[<大题数>]（默认 6，即「一」到「六」）
  \examscoretable[6]

  \vspace{0.6em}

  %% 大题与小题：
  %% \examsection{<栏目标题>}              自动加「一、」「二、」…（每题重新计小题号）
  %% \begin{examquestions}[<题间距>]        小题区，参数是题间距（默认 13pt），不是答题空间
  %%   \examquestion <题干>
  %%   \examanswerspace{7cm}                本题的答题空间（最后一题同样生效）
  %% \end{examquestions}
  %% \begin{examproblem} <题干> \end{examproblem}   单个大题
  %% \examfigure[<宽度>]{<图题>}{<图片文件>}        插图（图片放 figures/ 下）

  \examsection{〈第一栏标题，例如：填空题（本题共 ? 小题，每小题 ? 分，满分 ? 分）〉}
  \begin{examquestions}
    % 在这里写小题，例如：
    % \examquestion 若 $f(x)=x^2$，则 $f'(1)=$ \underline{\hspace{2cm}}。
  \end{examquestions}

  \examsection{〈第二栏标题，例如：计算题（本题共 ? 小题，每小题 ? 分，满分 ? 分）〉}
  \begin{examquestions}[0pt]%  0pt：题目紧排，答题空间逐题显式留出
    % \examquestion 〈题干〉
    % \examanswerspace{7cm}   本题 7cm 答题空间
  \end{examquestions}

  \examsection{〈第三栏标题〉}
  \begin{examquestions}
    % \examquestion 〈题干〉
    % \examfigure[0.6\textwidth]{〈图题〉}{figures/sample-figure.png}
  \end{examquestions}

\fi

\end{document}
